\begin{proof}[Proof of \cref{obj:thm:surrogate-certificate}]
Put
\[
A(q):=\sum_{k\in I_n}h_k(G_n)\{q_k^{-1}+(1-q_k)^{-1}\},
\qquad
C:=\max\{1,\epsilon^{-(\bar d-1)}\}.
\]
Then \(C\geq 1\), hence \(C\geq0\).  Recall also the three covariance loads, for \(i,j\in O_n\):
\[
r_{ij}^1(G_n,q)=\mathbf 1\{|S_{ij}(G_n)|>0\}\left(\prod_{k\in S_{ij}(G_n)}q_k^{-1}-1\right),
\quad
r_{ij}^0(G_n,q)=\mathbf 1\{|S_{ij}(G_n)|>0\}\left(\prod_{k\in S_{ij}(G_n)}(1-q_k)^{-1}-1\right),
\]
\[
r_{ij}^{10}(G_n)=\mathbf 1\{|S_{ij}(G_n)|>0\}.
\]

\begin{enumerate}
\item By \cref{obj:thm:convex-design}, whose positivity-margin and feasible-budget hypotheses are exactly those assumed here, the envelope-optimal design
\(p_n^*(G_n)\) is feasible and satisfies
\[
V_{\mathrm{env}}(G_n,p_n^*(G_n))
 \leq V_{\mathrm{env}}(G_n,q)
\quad\text{for every }q\in\mathcal P_{n,B_n,\epsilon}.
\]
% lean: CausalSmith.Experimentation.BipartiteMinimaxDesign.optimalDesign_feasible_minimizes

\item On \(\mathcal P_{n,B_n,\epsilon}\), every \(q_k\) and \(1-q_k\) lies in \([\epsilon,1-\epsilon]\), hence is strictly
positive.  Thus \(A\) is continuous there, being a finite sum of reciprocals of nonvanishing continuous
coordinate functions.  Nonemptiness and compactness from \cref{obj:thm:convex-design}
therefore yield a feasible surrogate design \(p_n^{\deg}(G_n)\) such that
\[
A(p_n^{\deg}(G_n))\leq A(q)
\quad\text{for every }q\in\mathcal P_{n,B_n,\epsilon}.
\]
% lean: CausalSmith.Experimentation.BipartiteMinimaxDesign.surrogateDesign_feasible_minimizes

\item Fix \(q\in\mathcal P_{n,B_n,\epsilon}\), and write
\[
K_{ij}:=r_{ij}^1(G_n,q)+r_{ij}^0(G_n,q)+2r_{ij}^{10}(G_n).
\]
Substituting the definition of \(h_k(G_n)\) into \(A(q)\) and exchanging the order of the finite summations,
\[
A(q)=\sum_{i\in O_n}\sum_{j\in O_n}\sum_{k\in S_{ij}(G_n)}
 n^{-1}|S_{ij}(G_n)|^{-1}
 \{q_k^{-1}+(1-q_k)^{-1}\},
\]
and \cref{obj:def:graph-envelope} gives
\[
\frac{V_{\mathrm{env}}(G_n,q)}4
=\sum_{i\in O_n}\sum_{j\in O_n}n^{-1}K_{ij}.
\]
For a fixed pair \((i,j)\), let \(S=S_{ij}(G_n)\).  If \(S=\varnothing\),
both corresponding terms are zero.  Otherwise, set
\(a_k=q_k^{-1}\) and \(b_k=(1-q_k)^{-1}\).  Since \(0<q_k<1\),
all these factors are at least one, so for each \(k\in S\) the product of the remaining factors is at least one and hence
\(a_k\le\prod_{l\in S}a_l\); averaging this over \(k\in S\) gives
\[
|S|^{-1}\sum_{k\in S}a_k\leq\prod_{k\in S}a_k,
\qquad
|S|^{-1}\sum_{k\in S}b_k\leq\prod_{k\in S}b_k,
\]
the second by the same argument applied to \((b_k)\).
For nonempty \(S\), the displayed definitions of the loads give
\[
K_{ij}=\left(\prod_{k\in S}a_k-1\right)+\left(\prod_{k\in S}b_k-1\right)+2
=\prod_{k\in S}a_k+\prod_{k\in S}b_k .
\]
Adding the preceding two inequalities yields
\[
|S|^{-1}\sum_{k\in S}(a_k+b_k)\leq K_{ij}.
\]
Multiplication by \(n^{-1}\geq0\) and summation over \((i,j)\) prove
\[
A(q)\leq \frac{V_{\mathrm{env}}(G_n,q)}4.
\]
% lean: CausalSmith.Experimentation.BipartiteMinimaxDesign.pairAverage_le_envelopeKernel

\item For the reverse bound, retain \(S\neq\varnothing\).  Since
\(S\subseteq N_i(G_n)\), \cref{obj:ass:bounded-outcome-degree}
gives \(|S|\leq d_i\le\bar d\).  Moreover, the feasibility bounds \(\epsilon\le q_k\le1-\epsilon\) give
\[
0\leq a_k\leq\epsilon^{-1},
\qquad
0\leq b_k\leq\epsilon^{-1} .
\]
After deleting any single index \(k\in S\), each of the two products has \(|S|-1\) factors and is therefore bounded by
\[
(\epsilon^{-1})^{|S|-1}
=\epsilon^{-(|S|-1)}
 \leq \epsilon^{-(\bar d-1)}
 \leq C,
\]
the first inequality because \(0<\epsilon\le1\) makes \(t\mapsto\epsilon^{-t}\) nondecreasing and \(|S|\le\bar d\).
For nonnegative factors whose product after removal of any one factor is at
most \(C\), the full product is at most \(C\) times the factors' average: for each \(k\in S\),
\(\prod_{l\in S}a_l=a_k\prod_{l\in S\setminus\{k\}}a_l\le Ca_k\), and averaging over \(k\in S\) gives the claim.
Applying this separately to \((a_k)\) and \((b_k)\) gives
\[
\prod_{k\in S}a_k
 \leq C|S|^{-1}\sum_{k\in S}a_k,
\qquad
\prod_{k\in S}b_k
 \leq C|S|^{-1}\sum_{k\in S}b_k.
\]
Adding these and using the identity for \(K_{ij}\) from step 3,
\[
K_{ij}\leq
C|S|^{-1}\sum_{k\in S}\{q_k^{-1}+(1-q_k)^{-1}\}.
\]
The same inequality is immediate when \(S=\varnothing\), both sides being zero.  Multiplying by
\(n^{-1}\geq0\), summing over \((i,j)\), and using the two displayed
expansions in the preceding step yields
\[
\frac{V_{\mathrm{env}}(G_n,q)}4\leq C\,A(q).
\]
This proves the asserted sandwich for every feasible \(q\).
% lean: CausalSmith.Experimentation.BipartiteMinimaxDesign.envelopeKernel_le_certificate_pairAverage

\item Let \(V_*:=V_{\mathrm{env}}(G_n,p_n^*(G_n))\).  If \(V_*>0\), the
two sandwich bounds and the minimizing property of \(p_n^{\deg}(G_n)\) from step 2 imply
\[
\frac{V_{\mathrm{env}}\!\left(G_n,p_n^{\deg}(G_n)\right)}4
 \leq C\,A(p_n^{\deg}(G_n))
 \leq C\,A(p_n^*(G_n))
 \leq C\,\frac{V_*}{4},
\]
the middle step because \(p_n^*(G_n)\) is feasible and \(C\ge0\).  Thus
\[
V_{\mathrm{env}}\!\left(G_n,p_n^{\deg}(G_n)\right)\leq C V_*,
\]
and division by \(V_*>0\), together with \cref{obj:def:approximation-ratio},
gives \(\alpha_{\mathrm{cert}}(G_n)\leq C\).  If \(V_*\not>0\), that
definition assigns \(\alpha_{\mathrm{cert}}(G_n)=1\), which is at most
\(C\).  The claimed approximation-ratio bound follows.
% lean: CausalSmith.Experimentation.BipartiteMinimaxDesign.surrogate_certificate
\end{enumerate}
\end{proof}
